% Generated draft: requires mathematical and source review.
\documentclass[11pt]{article}
\usepackage[margin=25mm]{geometry}
\usepackage{fontspec}
\setmainfont{Latin Modern Roman}
\usepackage{amsmath,amssymb,mathtools,unicode-math}
\setmathfont{Latin Modern Math}
\usepackage{enumitem,needspace,xurl,hyperref}
\setlength{\parindent}{0pt}
\setlength{\parskip}{5pt}
\setlength{\emergencystretch}{3em}
\newcommand{\N}{\mathbb{N}}
\newcommand{\Z}{\mathbb{Z}}
\newcommand{\Q}{\mathbb{Q}}
\newcommand{\R}{\mathbb{R}}
\newcommand{\C}{\mathbb{C}}
\newcommand{\F}{\mathbb{F}}
\newcommand{\A}{\mathbb{A}}
\newcommand{\PP}{\mathbb{P}}
\newcommand{\E}{\mathbb{E}}
\newcommand{\eps}{\varepsilon}
\newcommand{\dd}{\,\mathrm{d}}
\begin{document}
\section*{UnsolvedMath ID: 30005841}
\textbf{Source problem code:} catalogue entry
\subsection{Definitions and mathematical statement}
Let \(u:\Omega\to\R^n\) denote a deformation and let \(Du\) denote its deformation gradient. The source considers the elastic energy
\[
\mathcal E(u)=\int_\Omega W(x,Du(x))\,\dd x,
\]
where \(W\) is the stored-energy density. The dimension, hypotheses on the reference domain \(\Omega\), domain and range of \(W\), meaning of the derivative, and conditions ensuring that this integral is defined are not specified in the supplied record.

To express the variational question, write \(\mathcal A\) for an admissible class of deformations incorporating the function space, boundary conditions, and any physical constraints. This is a placeholder for missing problem data, not a class supplied by the source. For any specification of these data for which the energy and its infimum are meaningful, the requested existence property is
\[
\exists u_*\in\mathcal A\quad
\mathcal E(u_*)=\inf_{v\in\mathcal A}\mathcal E(v).
\]
The problem asks for conditions on the domain, stored-energy density, admissible class, and boundary data under which this infimum is attained, and under which an energy minimizer is an equilibrium solution of the corresponding equations of nonlinear elasticity.

The source does not specify an equilibrium equation, a notion of solution, loading, boundary conditions, or the regularity and admissible variations needed to connect minimization with equilibrium. It also does not specify whether a minimizer must have finite energy. Thus the displayed property is a schematic formulation of the existence question, rather than a fully specified conjecture. Regularity and the validity of equilibrium equations for minimizers require separate formulation and are not asserted to follow from attainment alone.

The supplied original excerpt places this question beside a contribution concerning nonlocal gradients in bounded domains and applications in finite and linear elasticity. It does not specify a nonlocal gradient or a nonlocal energy, so no such operator is substituted for \(Du\).
\subsection{Short English statement}
Determine assumptions under which the elastic energy has a minimizer among admissible deformations and under which that minimizer satisfies the intended nonlinear elasticity equilibrium equations. The source leaves the admissible deformations, boundary conditions, material hypotheses, and notion of equilibrium unspecified, so a unique precise theorem cannot yet be stated. Its supplied partially solved label is dated curation requiring independent review.
\subsection{Sources}
\begin{itemize}
\item Local UnsolvedMath registry: supplied record with permanent ID \(30005841\), title ``Equilibrium Solutions in Nonlinear Elastostatics,'' source problem code \texttt{catalogue entry}, in the Oberwolfach Reports Open Problems collection; proposed year: 2024. No individual problem proposer is supplied.
\item Supplied report citation: \emph{Fracture as an Emergent Phenomenon} (2024), Oberwolfach Reports, pp.\ 5--68, DOI 10.4171/owr/2024/1. \url{https://doi.org/10.4171/owr/2024/1}. The supplied background describes these reports as documenting mathematical workshops and their open-problem discussions. No specific page for the excerpt is supplied.
\item Attribution retained from the original excerpt: Jos\'e Carlos Bellido, ``Nonlocal gradients in bounded domains motivated by Continuum Mechanics: functional analysis framework and applications in finite and linear elasticity,'' joint work with J. Cueto and C. Mora-Corral. The record does not establish that these contributors proposed the extracted question.
\item Supplied literature source: ``Rate-form equilibrium for isotropic Cauchy elasticity.'' \url{https://doi.org/10.1007/s00332-025-10214-y}.
\item Supplied literature source: ``Some open problems in elasticity.'' \url{https://ora.ox.ac.uk/objects/uuid%3A4ae628ee-055d-47e3-bea0-888bc4063f40}.
\item Literature assessment dated exactly \texttt{2026-08-22}: the supplied curation labels the problem partially solved and reports existence of elastic-energy minimizers under coercivity and polyconvexity hypotheses, while describing the unrestricted nonlinear elastostatic existence, regularity, and equilibrium-equation problem as open. This is supplied dated curation derived from the OWR record, not an independently verified literature assessment; the status and cited claims require independent verification. No complete hypotheses for the reported existence results are supplied.
\item The registry labels the statement \texttt{corrected\_verified}, with the qualification that only formula delimiters, notation, grammar, and typography were repaired and that mathematical content and scope were preserved. This records the supplied editorial assessment, not mathematical verification of this draft.
\end{itemize}
\end{document}
